\section{The tiered access path: from collection to admitted segment}
\label{sec:retrieval}
\label{sec:lanes}

A query reaches evidence by narrowing its scope. Catalogue-first search
selects a collection, examines the descriptions on its shelf and opens
selected documents or passages. Pinned search starts within specified
collections and uses their indices to propose candidates. Both paths can
call the server-side Laya evidence gate. This design gives an agent
explicit choices at collection, document and passage granularity and
preserves the source identifiers needed to inspect those choices.

Figure~\ref{fig:tiered-algorithm} and Algorithm~\ref{alg:tiered} show how
an access policy selects, deduplicates and reads candidates, with or without
the evidence gate. The resulting pack is an input to answer assessment.

Access policies choose which operations to run. The intended workflow distinguishes
a completed review without support from an incomplete run with possible
partial evidence. This distinction depends on observable failure status;
the global-search defect in Section~\ref{sec:e7} shows it is not guaranteed
when a collection is silently skipped.

\begin{figure*}[t]
\centering
\includegraphics[width=\textwidth]{fig_tiered_drawio_20261116}
\caption{Tiered access workflow. Candidate references come from pinned
collections, catalogue navigation or a combination. The optional Laya gate
fetches and scores passages. A completed review without support reports
the searched scope. An incomplete operation reports failures and keeps
any available evidence.}
\label{fig:tiered-algorithm}
\end{figure*}

\begin{algorithm}[t]
\caption{Tiered access with optional evidence gate}
\label{alg:tiered}
\begin{algorithmic}[1]
\Require question $q$; catalogue $\mathcal{C}$; threshold $\tau$; gate flag $g$
\State $D \gets \textsc{GatherCandidates}(q,\mathcal{C})$ \Comment{pinned or catalogue-first}
\If{$g$}
  \State $S \gets \textsc{LayaJudge}(D,q,\tau)$ \Comment{fetch, score and return survivors}
\Else
  \State $S \gets \textsc{ReadSelected}(D)$
\EndIf
\If{execution incomplete}
  \State \Return incomplete status, searched scope, and available evidence
\ElsIf{completed review finds no support}
  \State \Return scoped not-found (searched scope, blind spots)
\Else
  \State \Return agent-synthesised answer with source references
\EndIf
\end{algorithmic}
\end{algorithm}
The policy instructs the agent to report its searched scope and remaining
blind spots when it reaches its call budget. The chat endpoint controls
timeouts and turn limits. Workflow instructions specify the retrieval
sequence and batch sizes.

\subsection{Tier policies and the evidence gate}
\label{sec:trace}
\label{sec:gate}
\label{sec:gate-primitive}
\label{sec:twostage}

The knowledge-base chat interface defines two policies. The search policy
starts from pinned collections and calls two-stage or vector search. The
librarian policy starts from the all-library catalogue, uses shelf
descriptions to choose document references and directs the agent to score
them in sequential batches before reading the survivors. The two-stage
tool uses lexical search and optional graph expansion to form a candidate
pool, applies dense ranking and returns candidates for inspection.

The MCP server also registers \path{kb_hybrid_search}, which merges vector
and catalogue candidates, optionally applies Laya, and sorts retained
documents by verifier score when judging is enabled. The librarian policy
directs the agent to use the catalogue and shelf tools directly. A separate
controlled hybrid experiment used per-library vector top-30 at cosine
threshold 0.35, catalogue candidates, fusion and a cap of twelve. The
ordering analysis in Section~\ref{sec:e6} replays historical hybrid pools
to compare ways of using the recorded scores.

The local Laya classifier \cite{laya}, exposed as \path{kb_laya_judge},
segments candidate text. Its \texttt{noul} output is a passage score on a
0--1 scale; passages at or above the threshold (0.5 by default) are eligible
for admission. It is not a calibrated probability that a passage supports
an answer claim. The gate groups survivors by source and position,
deduplicates them and applies a character cap. Where available, survivors
carry collection and document identifiers, file path, part or section, and
line range in the parsed text. These locators are not bound to immutable
versions. Fetch or scoring errors are reported separately from low scores.
Calling workflows may apply their own ordering. Table~\ref{tab:lanes}
summarises the policy choices.

\begin{table*}[t]
\centering
\small
\setlength{\tabcolsep}{5pt}
\caption{Access policies over the tiered path. Modes differ in how candidate collections and documents are selected.
The hybrid row describes the experimental composition used in the evaluation.}
\label{tab:lanes}
\begin{tabular}{@{}>{\raggedright\arraybackslash}p{0.11\textwidth}>{\raggedright\arraybackslash}p{0.19\textwidth}>{\raggedright\arraybackslash}p{0.24\textwidth}>{\raggedright\arraybackslash}p{0.16\textwidth}>{\raggedright\arraybackslash}p{0.20\textwidth}@{}}
\toprule
Policy & Tiers traversed & Retrieval affordance & Budget type & Intended terminal report\\
\midrule
Search & pinned collections, segment & Two-stage or vector search & Endpoint turn limit & Answer or scoped not-found\\
Librarian & catalogue, shelf, document, segment & Catalogue navigation and passage scoring & Workflow instructions & Answer or searched scope\\
Hybrid & catalogue and segment & Combined vector and catalogue candidates & Configuration-specific candidate cap & Candidates or judged survivors\\
\bottomrule
\end{tabular}
\end{table*}

Response transport also determines which candidates the agent can see.
The catalogue-truncation case exposed an incomplete preview before the agent
could inspect the whole shelf. The shelf tool now caps its response and
supplies grouped and paginated views, allowing the caller to request the
remaining entries. Section~\ref{sec:eval} reports the repair check.

\subsection{Honest abstention}
The answer workflow asks the agent to report absence after a completed
review that finds no support, including the scope it searched. A plausible
candidate may still lack the required evidence. If a tool is unavailable
or the search times out, the workflow reports the incomplete execution
and any usable evidence. Section~\ref{sec:eval} reports how the agents
handled these terminal conditions in the observed trials.
